% front.tex (국문 v13): 제목 · 저자 · 초록 · 핵심어(블록 1), 자료 가용성(블록 3), 뒷부분(블록 4).
\providecommand{\frontblock}{0}
\ifnum\frontblock=1
\title[물리 잔차 기반 기계학습의 연도별 활동층 두께 예측]{물리 잔차 기반 기계학습 워크플로를 활용한 영구동토 활동층 두께 예측과 지역별 편차 분석}

\author[1]{\fnm{백승원}}
\author[2]{\fnm{김승희}}
\author[2]{\fnm{정윤택}}
\author*[1]{\fnm{조용채}}\email{yc.cho@snu.ac.kr}
\affil*[1]{\orgdiv{에너지시스템공학부}, \orgname{서울대학교}, \orgaddress{\city{서울}, \postcode{08826}, \country{대한민국}}}
\affil[2]{\orgdiv{원격탐사빙권정보센터(Center of Remote Sensing and GIS)}, \orgname{극지연구소}, \orgaddress{\city{인천}, \postcode{21990}, \country{대한민국}}}

\abstract{활동층 두께(ALT) 관측은 드물고 공간적으로 편중되어 있어 지역 지도는 관측을 확장하는 모델 위에 서며, 흔히 쓰는 두 모델은 상충 관계에 있다. Stefan 모델은 관측이 없는 지역으로 전이되지만 토양을 계수 하나로 나타내고, 기계학습(ML)은 학습 지점을 잘 맞추지만 그 지점에서 떨어진 곳에서는 정확도가 낮아지며, ML 로 만든 공개 격자 제품도 같은 한계를 지닌다. 이 연구는 재보정한 Stefan 모델의 잔차를 기후 · 지형 · 토양 입력으로부터 gradient boosting 이 학습하는 물리 잔차 워크플로를, 알래스카 · 레나 델타 · 캐나다 북서부의 2010--2024년 지점-연도 관측 \nncyAll{}개에서 각 관측 연도의 기후를 써서 Stefan 모델, 직접 ML, 공개 제품 5종과 비교하였다. 공간 블록 검증에서 물리 잔차 모델은 Stefan 오차를 레나 델타에서 \ndrefLeblkRoabs{} cm, 캐나다에서 \ndrefCablkRoabs{} cm 줄였다. 한 지역을 학습에서 빼면 직접 ML 은 레나 델타에서 Stefan 오차를 \ndrefLeregDzabs{} cm 넘었으나 물리 잔차 모델은 Stefan 모델과 \ndrefLeregRoabs{} cm 안에 머물렀고, 지역 라벨 25개로 오차를 \ntlgainLeRstwfive{} cm 줄였으며, 라벨이 적을 때는 지역 안에서도 직접 ML 보다 정확하였다. 같은 지점-연도에서 물리 잔차 모델은 ML 제품보다 관측에 가까웠고, 레나 델타에서 그 차이가 가장 컸다. 이 워크플로는 적용 가능 영역과 예측 구간이 딸린 연도별 1 km 지도를 내놓으며, 관측이 적은 지역에 물리에 고정된 출발점을 주어 지역 라벨이 더해질수록 지도가 개선되게 한다.}

\keywords{영구동토, 활동층 두께, Stefan 모델, 물리 기반 기계학습, 공간 교차검증, 공간 전이성}
\fi

\ifnum\frontblock=3
\bmhead{자료 가용성}
모든 입력 자료는 공개 자료이다. ALT 관측: ABoVE 토양 수분 · 활동층 두께 v2(ORNL DAAC, \url{https://doi.org/10.3334/ORNLDAAC/2369}), ALLena 융해 깊이 자료(PANGAEA, \url{https://doi.org/10.1594/PANGAEA.973813}), GTN-P CALM(PANGAEA, \url{https://doi.org/10.1594/PANGAEA.972777}). 기후: ERA5-Land 월평균(Copernicus Climate Data Store). 공변량: Copernicus DEM GLO-30, SoilGrids 2.0, ESA CCI 영구동토 확률, MODIS MOD13Q1 · MOD10A2 v6.1(AppEEARS). 제품: Wei 등(\url{https://doi.org/10.5281/zenodo.21667583}), Liu 등(\url{https://doi.org/10.5281/zenodo.10070610}), ESA CCI ALT v4 · v5(CEDA), Aalto 등(\url{https://zenodo.org/records/5003804}), Yi와 Kimball(\url{https://doi.org/10.3334/ORNLDAAC/1760}). 지점-연도 라벨 표, 연도별 기후 입력 표, 정적 공변량 표, 모든 방법의 겹 밖 예측, 알래스카 · 레나 델타 · 캐나다 북서부의 연도별 1 km 지도(2011--2024년)는 게재 시 공개 저장소에 둔다[저장소 · DOI].
\fi

\ifnum\frontblock=4
\bmhead{감사의 글}
[연구비 · 감사의 글.]

\bmhead{저자 기여}
백승원은 연구를 설계하고 자료를 구축하였으며 분석을 수행하고 원고를 작성하였다. 김승희와 정윤택은 원격탐사 공변량과 영구동토 현장 자료의 해석에 기여하고 원고를 검토하였다. 조용채는 연구를 지도하고 원고를 검토하였다. 모든 저자가 최종 원고를 승인하였다.

\bmhead{추가 정보}
\textbf{보충 자료} 보충 방법 S1--S5, 보충 표 S1--S7, 보충 그림 S1--S11을 함께 제공한다.

\textbf{이해 상충} 저자들은 이해 상충이 없음을 밝힌다.
\fi
